\clearpage
\SourceLeaf{344}
\pdfbookmark[2]{7. Allgemeine Theorie der Berechnung der Cometenbahnen}{parabolisch-007}
\begin{center}
[7.]\par\medskip
Allgemeine Theorie der Berechnung der Cometenbahnen\par
[nebst Anwendung auf den zweiten Cometen des Jahrs 1813].
\end{center}
\begin{center}\small
\setlength{\tabcolsep}{3pt}
\begin{tabular}{c|c|c}
Beobachtungszeiten & Beobachtete Längen & Beobachtete Breiten\\\hline
\(t=[1813\ \text{April}]\ 7{,}55002\)&\(\alpha=271^{0}\,16'\,38''\)&\(\beta=+29^{0}\,2'\,0''\)\\
\(t'=\text{»}\quad\text{»}\quad14{,}54694\)&\(\alpha'=266\quad27\quad22\)&\(\beta'=+22\quad52\quad18\)\\
\(t''=\text{»}\quad\text{»}\quad21{,}59931\)&\(\alpha''=256\quad48\quad8\)&\(\beta''=+\ 9\quad53\quad12\)
\end{tabular}
\end{center}
\begin{center}
\begin{tabular}{c|c}
\shortstack{Längen der Sonne}&\shortstack{Logarithmen der Abstände\\von der Erde}\\\hline
\(\odot=17^{0}\,47'\,41''\)&\(\log R=0{,}00091\)\\
\(\odot'=24\quad38\quad45\)&\(\log R'=0{,}00175\)\\
\(\odot''=31\quad31\quad25\)&\(\log R''=0{,}00260\)
\end{tabular}
\end{center}
\noindent[\(r,r',r''\) die Abstände von der Sonne,\par
\noindent\(\rho,\rho',\rho''\) die Abstände von der Erde.]
\begin{center}Vorläufige Operationen:\end{center}
[Es ist
\[
\begin{aligned}
rr&=\rho\rho-2\rho R\cos\beta\cos(\alpha-\odot)+RR\\
r''r''&=\rho''\rho''-2\rho''R''\cos\beta''\cos(\alpha''-\odot'')+R''R'';
\end{aligned}
\]
setzt man, wie in der Abhandlung »Observationes cometae secundi A. MDCCCXIII« etc.]
\[
\begin{array}{l@{\quad}l}
[\mathrm{I}]&\begin{aligned}
\cos\psi&=\cos\beta\cos(\alpha-\odot),& A&=R\sin\psi\\
\cos\psi''&=\cos\beta''\cos(\alpha''-\odot''),& A''&=R''\sin\psi'',
\end{aligned}
\end{array}
\]
[so wird
\[
\begin{aligned}
rr&=(\rho-R\cos\psi)^2+AA\\
r''r''&=(\rho''-R''\cos\psi'')^2+A''A''.
\end{aligned}
\]

Legt man durch den ersten und dritten geocentrischen Ort einen grössten Kreis, der die Ekliptik im Punkte \(N\) unter dem Neigungswinkel \(J\) schneide, so ist, wenn die Länge dieses Punkts ebenfalls mit \(N\) bezeichnet wird,
\SourceLeaf{345}
\[
\begin{aligned}
\operatorname{tang}J\sin(\alpha-N)&=\operatorname{tang}\beta\\
\operatorname{tang}J\sin(\alpha''-N)&=\operatorname{tang}\beta'';
\end{aligned}
\]
\(J\) und \(N\) können aus den Formeln, vgl. Theoria motus corporum coelestium Art. 78,]
\[
\begin{array}{l@{\quad}l}
[\mathrm{II}]&\begin{aligned}
\operatorname{tang}J\sin(\alpha-N)&=\operatorname{tang}\beta\\
\operatorname{tang}J\cos(\alpha-N)&=\frac{\operatorname{tang}\beta''-\operatorname{tang}\beta\cos(\alpha''-\alpha)}{\sin(\alpha''-\alpha)}
\end{aligned}
\end{array}
\]
[berechnet werden.

Indem man die Längen vom Punkte \(N\) aus zählt, hat man für die Chorde zwischen dem ersten und dritten Ort
\[
\begin{aligned}
\text{Chorde}^2={}&\bigl\{\rho''\cos\beta''\cos(\alpha''-N)-\rho\cos\beta\cos(\alpha-N)\\
&\qquad{}-R''\cos(\odot''-N)+R\cos(\odot-N)\bigr\}^2\\
&+\bigl\{\rho''\cos\beta''\sin(\alpha''-N)-\rho\cos\beta\sin(\alpha-N)\\
&\qquad{}-R''\sin(\odot''-N)+R\sin(\odot-N)\bigr\}^2\\
&+\{\rho''\sin\beta''-\rho\sin\beta\}^2.
\end{aligned}
\]

Man setze
\[
\begin{aligned}
R''\cos(\odot''-N)-R\cos(\odot-N)&=f\cos(F-N)\\
R''\sin(\odot''-N)-R\sin(\odot-N)&=f\sin(F-N)
\end{aligned}
\]
und bezeichne mit \(K,K''\) die Abstände des ersten und dritten geocentrischen Orts vom Punkte \(N\); alsdann ist
\[
\begin{aligned}
\cos\beta\cos(\alpha-N)&=\cos K,\\
\cos\beta\sin(\alpha-N)&=\cos J\sin K,\\
\sin\beta&=\sin J\sin K\\
\cos\beta''\cos(\alpha''-N)&=\cos K'',\\
\cos\beta''\sin(\alpha''-N)&=\cos J\sin K'',\\
\sin\beta''&=\sin J\sin K'',
\end{aligned}
\]
womit wird
\[
\begin{aligned}
\text{Chorde}^2={}&(\rho''\cos K''-\rho\cos K-f\cos(F-N))^2\\
&+(\rho''\cos J\sin K''-\rho\cos J\sin K-f\sin(F-N))^2\\
&+(\rho''\sin J\sin K''-\rho\sin J\sin K)^2.
\end{aligned}
\]

Man kann die Grössen \(f,F,K,K''\) aus den Formeln]
\[
\begin{array}{l@{\quad}l}
[\mathrm{III}]&\begin{aligned}
R''\cos(\odot''-\odot)-R&=f\cos(F-\odot)\\
R''\sin(\odot''-\odot)&=f\sin(F-\odot)
\end{aligned}
\end{array}
\]
\[
\begin{array}{l@{\quad}l}
[\mathrm{IV}]&\begin{aligned}
\frac{\operatorname{tang}(\alpha-N)}{\cos J}&=\operatorname{tang}K\\
\frac{\operatorname{tang}(\alpha''-N)}{\cos J}&=\operatorname{tang}K''
\end{aligned}
\end{array}
\]
\SourceLeaf{346}
[berechnen und noch setzen]
\[
\begin{array}{l@{\quad}l}
[\mathrm{V}]&\begin{aligned}
f\cos J\sin(F-N)&=g\sin G\\
f\cos(F-N)&=g\cos G.
\end{aligned}
\end{array}
\]

[Legt man nun durch den mittlern geocentrischen Ort einen grössten Kreis, welcher den durch die beiden äussern Örter gehenden grössten Kreis rechtwinklig schneidet und nennt man \(Q\) die Länge des Punkts, in dem dieser Kreis die Ekliptik schneidet, so ist]
\[
\begin{array}{l@{\quad}l}
[\mathrm{VI}]&\begin{aligned}
\operatorname{tang}(Q-N)&=\operatorname{tang}(\alpha'-N)\cdot\left(1+\frac{\operatorname{tang}J\operatorname{tang}\beta'}{\sin(\alpha'-N)}\right)\\
&=\operatorname{tang}(\alpha'-N)+\frac{\operatorname{tang}J\operatorname{tang}\beta'}{\cos(\alpha'-N)}\mathpunct{\cdot}
\end{aligned}
\end{array}
\]

[In unserm Beispiel wird]
\[
\begin{array}{l@{\;}l@{\;}r}
\cos\beta&\mathord{.}\,\mathord{.}\,\mathord{.}\,\mathord{.}\,\mathord{.}\,\mathord{.}\,\mathord{.}\,\mathord{.}\,\mathord{.}\,\mathord{.}&9{,}94168\\
\cos(\alpha-\odot)&\mathord{.}\,\mathord{.}\,\mathord{.}\,\mathord{.}&9{,}45379_n\\\hline
\cos\psi&\mathord{.}\,\mathord{.}\,\mathord{.}\,\mathord{.}\,\mathord{.}\,\mathord{.}\,\mathord{.}\,\mathord{.}\,\mathord{.}&9{,}39547_n\\
R&\mathord{.}\,\mathord{.}\,\mathord{.}\,\mathord{.}\,\mathord{.}\,\mathord{.}\,\mathord{.}\,\mathord{.}\,\mathord{.}\,\mathord{.}\,\mathord{.}\,\mathord{.}&0{,}00091\\
\sin\psi&\mathord{.}\,\mathord{.}\,\mathord{.}\,\mathord{.}\,\mathord{.}\,\mathord{.}\,\mathord{.}\,\mathord{.}\,\mathord{.}\,\mathord{.}&9{,}98615\\\hline
A&\mathord{.}\,\mathord{.}\,\mathord{.}\,\mathord{.}\,\mathord{.}\,\mathord{.}\,\mathord{.}\,\mathord{.}\,\mathord{.}\,\mathord{.}\,\mathord{.}\,\mathord{.}&9{,}98706
\end{array}
\]
\[
\begin{array}{l@{\;}l@{\;}r}
\cos\beta''&\mathord{.}\,\mathord{.}\,\mathord{.}\,\mathord{.}\,\mathord{.}\,\mathord{.}\,\mathord{.}\,\mathord{.}\,\mathord{.}&9{,}99350\\
\cos(\alpha''-\odot'')&\mathord{.}\,\mathord{.}\,\mathord{.}&9{,}84736_n\\\hline
\cos\psi''&\mathord{.}\,\mathord{.}\,\mathord{.}\,\mathord{.}\,\mathord{.}\,\mathord{.}\,\mathord{.}\,\mathord{.}\,\mathord{.}&9{,}84086_n\\
R''&\mathord{.}\,\mathord{.}\,\mathord{.}\,\mathord{.}\,\mathord{.}\,\mathord{.}\,\mathord{.}\,\mathord{.}\,\mathord{.}\,\mathord{.}\,\mathord{.}&0{,}00260\\
\sin\psi''&\mathord{.}\,\mathord{.}\,\mathord{.}\,\mathord{.}\,\mathord{.}\,\mathord{.}\,\mathord{.}\,\mathord{.}\,\mathord{.}&9{,}85778\\\hline
A''&\mathord{.}\,\mathord{.}\,\mathord{.}\,\mathord{.}\,\mathord{.}\,\mathord{.}\,\mathord{.}\,\mathord{.}\,\mathord{.}\,\mathord{.}\,\mathord{.}&9{,}86038
\end{array}
\]
\[
\begin{array}{l@{\;}l@{\;}r}
R''&\mathord{.}\,\mathord{.}\,\mathord{.}\,\mathord{.}\,\mathord{.}\,\mathord{.}\,\mathord{.}\,\mathord{.}\,\mathord{.}\,\mathord{.}\,\mathord{.}&0{,}00260\\
\cos(\odot''-\odot)&\mathord{.}\,\mathord{.}\,\mathord{.}&9{,}98741\\
\sin(\odot''-\odot)&\mathord{.}\,\mathord{.}\,\mathord{.}&9{,}37535\\\hline
[R''\cos(\odot''-\odot)]&&9{,}99001\\
-R&\mathord{.}\,\mathord{.}\,\mathord{.}\,\mathord{.}\,\mathord{.}\,\mathord{.}\,\mathord{.}\,\mathord{.}\,\mathord{.}\,\mathord{.}&0{,}00091_n\\\hline
\left[\begin{array}{@{}l@{}}f\cos(F-\odot)\\f\sin(F-\odot)\end{array}\right]&
\begin{array}{@{}l@{}}\mathord{.}\\\mathord{.}\end{array}&
\begin{array}{@{}r@{}}8{,}39501_n\\9{,}37795\end{array}\\\hline
\sin(F-\odot)&\mathord{.}\,\mathord{.}\,\mathord{.}\,\mathord{.}&9{,}99766\\
f&\mathord{.}\,\mathord{.}\,\mathord{.}\,\mathord{.}\,\mathord{.}\,\mathord{.}\,\mathord{.}\,\mathord{.}\,\mathord{.}\,\mathord{.}\,\mathord{.}\,\mathord{.}\,\mathord{.}&9{,}38029\\
\multicolumn{2}{r}{F-\odot=}&95^{0}\,56'\,16''\\
\multicolumn{2}{r}{\odot=}&17\quad47\quad41\\\cline{2-3}
\multicolumn{2}{r}{F=}&113\quad43\quad57
\end{array}
\]
\[
\begin{array}{l@{\;}l@{\;}r}
\operatorname{tang}\beta&\mathord{.}\,\mathord{.}\,\mathord{.}\,\mathord{.}\,\mathord{.}\,\mathord{.}\,\mathord{.}\,\mathord{.}\,\mathord{.}\,\mathord{.}\,\mathord{.}&9{,}74435\\
\cos(\alpha''-\alpha)&\mathord{.}\,\mathord{.}\,\mathord{.}\,\mathord{.}\,\mathord{.}\,\mathord{.}\,\mathord{.}&9{,}98599\\\hline
&&9{,}73034_n\\
\operatorname{tang}\beta''&\mathord{.}\,\mathord{.}\,\mathord{.}\,\mathord{.}\,\mathord{.}\,\mathord{.}\,\mathord{.}\,\mathord{.}\,\mathord{.}\,\mathord{.}&9{,}24127\\\hline
&&9{,}56010_n\\
\sin(\alpha''-\alpha)&\mathord{.}\,\mathord{.}\,\mathord{.}\,\mathord{.}\,\mathord{.}\,\mathord{.}\,\mathord{.}&9{,}39786_n\\\hline
\left[\begin{array}{@{}l@{}}\operatorname{tang}J\cos(\alpha-N)\\\operatorname{tang}J\sin(\alpha-N)\end{array}\right]&&
\begin{array}{@{}r@{}}0{,}16224\\9{,}74435\end{array}\\\hline
\cos(\alpha-N)&\mathord{.}\,\mathord{.}\,\mathord{.}\,\mathord{.}\,\mathord{.}\,\mathord{.}\,\mathord{.}&9{,}97041\\
\operatorname{tang}J&\mathord{.}\,\mathord{.}\,\mathord{.}\,\mathord{.}\,\mathord{.}\,\mathord{.}\,\mathord{.}\,\mathord{.}\,\mathord{.}\,\mathord{.}&0{,}19183\\
\sin J&\mathord{.}\,\mathord{.}\,\mathord{.}\,\mathord{.}\,\mathord{.}\,\mathord{.}\,\mathord{.}\,\mathord{.}\,\mathord{.}\,\mathord{.}\,\mathord{.}\,\mathord{.}&9{,}92487\\
\cos J&\mathord{.}\,\mathord{.}\,\mathord{.}\,\mathord{.}\,\mathord{.}\,\mathord{.}\,\mathord{.}\,\mathord{.}\,\mathord{.}\,\mathord{.}\,\mathord{.}\,\mathord{.}&9{,}73304\\
\multicolumn{2}{r}{\alpha-N=}&20^{0}\,54'\,32''
\end{array}
\]
\[
\begin{array}{r@{\;}l}
N=&250^{0}\,22'\,6''\\
\alpha'-N=&16\quad5\quad16\\
\alpha''-N=&6\quad26\quad2\\
F-N=&223\quad21\quad51
\end{array}
\]
\SourceLeaf{347}
\[
\begin{array}{l@{\;}l@{\;}r}
\operatorname{tang}(\alpha-N)&\mathord{.}\,\mathord{.}\,\mathord{.}\,\mathord{.}\,\mathord{.}&9{,}58211\\
\operatorname{tang}(\alpha''-N)&\mathord{.}\,\mathord{.}\,\mathord{.}\,\mathord{.}&9{,}05218\\
\cos J&\mathord{.}\,\mathord{.}\,\mathord{.}\,\mathord{.}\,\mathord{.}\,\mathord{.}\,\mathord{.}\,\mathord{.}\,\mathord{.}\,\mathord{.}\,\mathord{.}&9{,}73304\\\hline
\multicolumn{2}{r}{K=}&35^{0}\,14'\,19''\\
\multicolumn{2}{r}{K''=}&11\quad46\quad42
\end{array}
\]
\[
\begin{array}{l@{\;}l@{\;}r}
\cos J&\mathord{.}\,\mathord{.}\,\mathord{.}\,\mathord{.}\,\mathord{.}\,\mathord{.}\,\mathord{.}\,\mathord{.}\,\mathord{.}\,\mathord{.}\,\mathord{.}&9{,}73304\\
f&\mathord{.}\,\mathord{.}\,\mathord{.}\,\mathord{.}\,\mathord{.}\,\mathord{.}\,\mathord{.}\,\mathord{.}\,\mathord{.}\,\mathord{.}\,\mathord{.}\,\mathord{.}\,\mathord{.}\,\mathord{.}&9{,}38029\\
\sin(F-N)&\mathord{.}\,\mathord{.}\,\mathord{.}\,\mathord{.}\,\mathord{.}&9{,}83672_n\\
\cos(F-N)&\mathord{.}\,\mathord{.}\,\mathord{.}\,\mathord{.}\,\mathord{.}&9{,}86154_n\\\hline
\left[\begin{array}{@{}l@{}}g\sin G\\g\cos G\end{array}\right]&
\begin{array}{@{}l@{}}\mathord{.}\,\mathord{.}\,\mathord{.}\,\mathord{.}\,\mathord{.}\,\mathord{.}\,\mathord{.}\\\mathord{.}\,\mathord{.}\,\mathord{.}\,\mathord{.}\,\mathord{.}\,\mathord{.}\,\mathord{.}\end{array}&
\begin{array}{@{}r@{}}8{,}95005_n\\9{,}24183_n\end{array}\\\hline
\cos G&\mathord{.}\,\mathord{.}\,\mathord{.}\quad\mathord{.}\,\mathord{.}\,\mathord{.}\,\mathord{.}\,\mathord{.}\,\mathord{.}&9{,}94966_n\\
g&\mathord{.}\,\mathord{.}\,\mathord{.}\,\mathord{.}\,\mathord{.}\,\mathord{.}\,\mathord{.}\,\mathord{.}\,\mathord{.}\,\mathord{.}\,\mathord{.}\,\mathord{.}\,\mathord{.}&9{,}29217\\
\multicolumn{2}{r}{G=}&207^{0}\,3'\,23''
\end{array}
\]
\[
\begin{array}{l@{\;}l@{\;}r}
\sin(\alpha'-N)&\mathord{,}\,\mathord{.}\,\mathord{.}\,\mathord{.}\,\mathord{.}\,\mathord{.}&9{,}44265\\
\operatorname{cotang}J&\mathord{.}\,\mathord{.}\,\mathord{.}\,\mathord{.}\,\mathord{.}\,\mathord{.}\,\mathord{.}\,\mathord{.}&9{,}80817\\
\operatorname{cotang}\beta'&\mathord{.}\,\mathord{.}\,\mathord{.}\,\mathord{.}\,\mathord{.}\,\mathord{.}\,\mathord{.}\,\mathord{.}&0{,}37486\\\hline
&&9{,}62568\\
\multicolumn{3}{l}{\left[1+\dfrac{\operatorname{tang}J\operatorname{tang}\beta'}{\sin(\alpha'-N)}\quad\mathord{.}\,\mathord{.}\quad0{,}52733\right]}\\
\operatorname{tang}(\alpha'-N)&\mathord{.}\,\mathord{.}\,\mathord{.}\,\mathord{.}&9{,}46000\\\hline
\operatorname{tang}(Q-N)&\mathord{.}\,\mathord{.}\,\mathord{.}\,\mathord{.}&9{,}98733
\end{array}
\]
\[
\begin{array}{r@{\;}l}
Q-N=&44^{0}\,9'\,51''\\
Q=&294\quad31\quad57\\
\odot'-Q=&90\quad6\quad48\\
\alpha-Q=&-23\quad15\quad19\\
\alpha'-Q=&-28\quad4\quad35\\
\alpha''-Q=&-37\quad43\quad49\\
\odot-Q=&83\quad15\quad44\\
\odot''-Q=&96\quad59\quad28.
\end{array}
\]

Die bisherigen Vorschriften sind allgemein. Jetzt muss man sich entschliessen, was man von den Beobachtungsdatis weglassen will. Man denke sich durch den mittelsten geocentrischen Ort zwei grösste Kreise gezogen. Er befindet sich also wirklich in beiden zugleich, allein da man zu viele Data hat, so wird man die eine Bedingung aufopfern, und bloss die beibehalten, dass er in \textls[150]{einem} gegebenen grössten Kreise liegt. Dieser ist an sich ganz willkürlich; legt man ihn durch den [mittelsten] Ort der Sonne, so hat man die bequemste Rechnung, aber zuweilen nicht hinlängliche Schärfe. Allgemein schneide derselbe die Ekliptik in dem Punkte, dessen Länge \(=P\). Setzt man \(P=N\), so wird bloss die Krümmung der scheinbaren Cometenbahn dargestellt, aber nicht die veränderte Geschwindigkeit. Setzt man \(P=Q\), so verhält sichs umgekehrt, und dies wird im allgemeinen das schärfste Resultat geben.

Wir haben nunmehro, wenn wir \(n,n',n'',k\) in der Bedeutung der Th[eoria] M[otus] C[orporum] C[oelestium] nehmen, hingegen \(\rho,\rho''\) die wahren Abstände des Cometen [von der Erde] bedeuten lassen [und wenn wir die Längen von dem noch ganz willkürlichen Punkte \(P\) zählen, nach Art. 112 der genannten
\SourceLeaf{348}Theoria
\[
\begin{aligned}
&n\rho\cos\beta\sin(\alpha-P)-n'\rho'\cos\beta'\sin(\alpha'-P)\\
&\qquad{}+n''\rho''\cos\beta''\sin(\alpha''-P)\\
&\quad=nR\sin(\odot-P)-n'R'\sin(\odot'-P)+n''R''\sin(\odot''-P)\\
&n\rho\sin\beta-n'\rho'\sin\beta'+n''\rho''\sin\beta''=0
\end{aligned}
\]
und hieraus zur Herstellung einer Beziehung zwischen \(\rho\) und \(\rho''\)] die vollkommen strenge Gleichung
\[
\begin{array}{l@{\quad}l}
[\mathrm{VII}]&\begin{aligned}
&n\rho\sin\beta\left(\frac{\sin(\alpha-P)}{\operatorname{tang}\beta}-\frac{\sin(\alpha'-P)}{\operatorname{tang}\beta'}\right)\\
&{}+n''\rho''\sin\beta''\left(\frac{\sin(\alpha''-P)}{\operatorname{tang}\beta''}-\frac{\sin(\alpha'-P)}{\operatorname{tang}\beta'}\right)\\
&=nR\sin(\odot-P)-n'R'\sin(\odot'-P)\\
&\qquad{}+n''R''\sin(\odot''-P).
\end{aligned}
\end{array}
\]

Nun kann man ohne Bedenken setzen [vgl. Art. 133 und 134 der Theoria motus]
\[
\begin{aligned}
\frac{n''}{n}&=\frac{t'-t}{t''-t'}\\
n'&=(n+n'')\left(1-\frac{4kk(t'-t)(t''-t')}{(r+r'')^3}\right)
\end{aligned}
\]
[oder
\[
\frac{n'}{n}=\frac{t''-t}{t''-t'}-\frac{4kk(t''-t)(t'-t)}{(r+r'')^3}\Biggr]\mathpunct{\cdot}
\]
Wir haben also
\[
\begin{array}{l@{\quad}l}
[\mathrm{VIII}]&\displaystyle\frac{1}{(r+r'')^3}=\mu+M\rho+M''\rho'',
\end{array}
\]
wenn wir setzen
\[
\begin{array}{l@{\quad}l}
[\mathrm{IX}]&\displaystyle M=\frac{\sin\beta\left(\dfrac{\sin(\alpha-P)}{\operatorname{tang}\beta}-\dfrac{\sin(\alpha'-P)}{\operatorname{tang}\beta'}\right)}{4kk(t''-t)(t'-t)R'\sin(\odot'-P)}
\end{array}
\]
\[
\begin{array}{l@{\quad}l}
[\mathrm{X}]&\displaystyle M''=\frac{\sin\beta''\left(\dfrac{\sin(\alpha''-P)}{\operatorname{tang}\beta''}-\dfrac{\sin(\alpha'-P)}{\operatorname{tang}\beta'}\right)}{4kk(t''-t)(t''-t')R'\sin(\odot'-P)}
\end{array}
\]
\[
\begin{array}{l@{\quad}l}
[\mathrm{XI}]&\displaystyle\mu=-\frac{\begin{gathered}(t''-t')R\sin(\odot-P)-(t''-t)R'\sin(\odot'-P)\\{}+(t'-t)R''\sin(\odot''-P)\end{gathered}}{4kk(t''-t)(t'-t)(t''-t')R'\sin(\odot'-P)}\mathpunct{\cdot}
\end{array}
\]

Man kann indessen hier füglich annehmen \(\mu=\dfrac{1}{(R+R'')^3}\) oder \(\mu=\dfrac{1}{8R'^3}\) [denn die Gleichungen VII und VIII gelten auch, die letztere genähert, wenn man anstatt aller auf den Cometen sich beziehender Data die entsprechenden für die Erde setzt, wobei \(\rho\) und \(\rho''\) verschwinden und der vorstehende Näherungswert sich ergibt.

Unser Beispiel gibt, wenn \(P=272^{0}\,40'\,21''\) genommen wird:]
\SourceLeaf{349}
\[
\begin{array}{l@{\;}l@{\;}r}
4kk&\mathord{.}\,\mathord{.}\,\mathord{.}\,\mathord{.}\,\mathord{.}\,\mathord{.}&7{,}07322\\
R'&\mathord{.}\,\mathord{.}\,\mathord{.}\,\mathord{.}\,\mathord{.}\,\mathord{.}\,\mathord{.}&0{,}00175\\
\sin(\odot'-P)&&9{,}96725\\\hline
&&7{,}04222\\
t''-t&\mathord{.}\,\mathord{.}\,\mathord{.}\,\mathord{.}\,\mathord{.}&1{,}14766\\\hline
&&8{,}18988\\
t''-t'&\mathord{.}\,\mathord{.}\,\mathord{.}\,\mathord{.}\,\mathord{.}&0{,}84833\\
t'-t&\mathord{.}\,\mathord{.}\,\mathord{.}\,\mathord{.}\,\mathord{.}&0{,}84491
\end{array}
\]
\[
\begin{array}{l@{\;}l@{\;}r}
\sin(\alpha-P)&&8{,}38650_n\\
\operatorname{tang}\beta&\mathord{.}\,\mathord{.}\,\mathord{.}&9{,}74435\\
\sin(\alpha'-P)&&9{,}03456_n\\
\operatorname{tang}\beta'&\mathord{.}\,\mathord{.}\,\mathord{.}&9{,}62514\\
\sin(\alpha''-P)&&9{,}43689_n\\
\operatorname{tang}\beta''&\mathord{.}\,\mathord{.}\,\mathord{.}&9{,}24127
\end{array}
\]
\[
\begin{array}{l@{\;}l@{\;}r}
\left[\dfrac{\sin(\alpha-P)}{\operatorname{tang}\beta}\right]&\mathord{.}\,\mathord{.}\,\mathord{.}\,\mathord{.}&8{,}64215_n\\[3pt]
\left[-\dfrac{\sin(\alpha'-P)}{\operatorname{tang}\beta'}\right]&\mathord{.}\,\mathord{.}&9{,}40942\\[3pt]\hline
&&9{,}32803\\
\sin\beta&\mathord{.}\,\mathord{.}\,\mathord{.}\,\mathord{.}\,\mathord{.}\,\mathord{.}\,\mathord{.}\,\mathord{.}\,\mathord{.}&9{,}68603\\\hline
&&9{,}01406\\
{}[\text{Nenner}]&\mathord{.}\,\mathord{.}\,\mathord{.}\,\mathord{.}\,\mathord{.}&9{,}03479\\\hline
M&\mathord{.}\,\mathord{.}\,\mathord{.}\,\mathord{.}\,\mathord{.}\,\mathord{.}\,\mathord{.}\,\mathord{.}\,\mathord{.}\,\mathord{.}&9{,}97927
\end{array}
\]
\[
\begin{array}{l@{\;}l@{\;}r}
\left[\dfrac{\sin(\alpha''-P)}{\operatorname{tang}\beta''}\right]&\mathord{.}\,\mathord{.}\,\mathord{.}\,\mathord{.}&0{,}19562_n\\[3pt]
&\mathord{.}\,\mathord{.}\,\mathord{.}\,\mathord{.}\,\mathord{.}\,\mathord{.}\,\mathord{.}\,\mathord{.}\,\mathord{.}\,\mathord{.}\,\mathord{.}\,\mathord{.}&9{,}40942\\\hline
&&0{,}11803_n\\
\sin\beta''&\mathord{.}\,\mathord{.}\,\mathord{.}\,\mathord{.}\,\mathord{.}\,\mathord{.}\,\mathord{.}\,\mathord{.}&9{,}23477\\\hline
&&9{,}35280_n\\
{}[\text{Nenner}]&\mathord{.}\,\mathord{.}\,\mathord{.}\,\mathord{.}\,\mathord{.}&9{,}03821\\\hline
M''&\mathord{.}\,\mathord{.}\,\mathord{.}\,\mathord{.}\,\mathord{.}\,\mathord{.}\,\mathord{.}\,\mathord{.}\,\mathord{.}\,\mathord{.}&0{,}31459_n
\end{array}
\]
\[
\begin{array}{l@{\;}l@{\;}r}
\sin(\odot-P)&\mathord{.}\,\mathord{.}&9{,}98470\\
-R&\mathord{.}\,\mathord{.}\,\mathord{.}\,\mathord{.}\,\mathord{.}\,\mathord{.}\,\mathord{.}\,\mathord{.}&0{,}00091_n\\
\text{C log Denom.}&\mathord{.}&0{,}96521\\\hline
&&0{,}95082_n
\end{array}
\]
\[
\begin{array}{l@{\;}l@{\;}r}
\sin(\odot'-P)&\mathord{.}&9{,}96725\\
R'&\mathord{.}\,\mathord{.}\,\mathord{.}\,\mathord{.}\,\mathord{.}\,\mathord{.}\,\mathord{.}\,\mathord{.}\,\mathord{.}&0{,}00175\\
&\mathord{.}\,\mathord{.}\,\mathord{.}\,\mathord{.}\,\mathord{.}\,\mathord{.}\,\mathord{.}\,\mathord{.}\,\mathord{.}\,\mathord{.}&1{,}26454\\\hline
&&1{,}23354
\end{array}
\]
\[
\begin{array}{l@{\;}l@{\;}r}
\sin(\odot''-P)&\mathord{.}\,\mathord{.}&9{,}94245\\
-R''&\mathord{.}\,\mathord{.}\,\mathord{.}\,\mathord{.}\,\mathord{.}\,\mathord{.}\,\mathord{.}\,\mathord{.}&0{,}00260_n\\
&\mathord{.}\,\mathord{.}\,\mathord{.}\,\mathord{.}\,\mathord{.}\,\mathord{.}\,\mathord{.}\,\mathord{.}\,\mathord{.}\,\mathord{.}\,\mathord{.}&0{,}96179\\\hline
&&0{,}90684_n
\end{array}
\]
\[
\begin{array}{r@{\;}l@{\;}r}
\mu&\mathord{.}\,\mathord{.}\,\mathord{.}\,\mathord{.}&9{,}0888\\
\text{aus der Näherungsformel: }\mu&\mathord{.}\,\mathord{.}\,\mathord{.}\,\mathord{.}&9{,}09166.
\end{array}
\]

Die übrigen Vorbereitungsrechnungen sind folgende:

[Man setze
\[
\begin{aligned}
M\sin K''+M''\sin K&=h\sin H\\
M\cos K''+M''\cos K&=h\cos H,
\end{aligned}
\]
oder für die Rechnung bequemer]
\[
\begin{array}{l@{\quad}l}
[\mathrm{XII}]&\begin{aligned}
M\sin(K''-K)&=h\sin(H-K)\\
M\cos(K''-K)+M''&=h\cos(H-K)
\end{aligned}
\end{array}
\]
\[
\begin{array}{l@{\;}l@{\;}r}
M&\mathord{.}\,\mathord{.}\,\mathord{.}\,\mathord{.}\,\mathord{.}\,\mathord{.}\,\mathord{.}\,\mathord{.}\,\mathord{.}\,\mathord{.}\,\mathord{.}\,\mathord{.}\,\mathord{.}&9{,}97927\\
\sin(K''-K)&\mathord{.}\,\mathord{.}\,\mathord{.}\,\mathord{.}\,\mathord{.}&9{,}60001_n\\
\cos(K''-K)&\mathord{.}\,\mathord{.}\,\mathord{.}\,\mathord{.}\,\mathord{.}&9{,}96253\\\hline
\left[\begin{array}{@{}l@{}}M\cos(K''-K)\\M''\quad\mathord{.}\,\mathord{.}\,\mathord{.}\,\mathord{.}\,\mathord{.}\,\mathord{.}\,\mathord{.}\,\mathord{.}\,\mathord{.}\,\mathord{.}\end{array}\right]&
\begin{array}{@{}l@{}}\mathord{.}\\\mathord{.}\end{array}&
\begin{array}{@{}r@{}}9{,}94180\\0{,}31459_n\end{array}
\end{array}
\]
\[
\begin{array}{l@{\;}l@{\;}r}
\left[\begin{array}{@{}l@{}}h\sin(H-K)\\h\cos(H-K)\end{array}\right]&
\begin{array}{@{}l@{}}\mathord{.}\,\mathord{.}\\\mathord{.}\,\mathord{.}\end{array}&
\begin{array}{@{}r@{}}9{,}57928_n\\0{,}07513_n\end{array}\\\hline
[\cos(H-K)]&\mathord{.}\,\mathord{.}\,\mathord{.}\,\mathord{.}&9{,}97892_n\\
h&\mathord{.}\,\mathord{.}\,\mathord{.}\,\mathord{.}\,\mathord{.}\,\mathord{.}\,\mathord{.}\,\mathord{.}\,\mathord{.}\,\mathord{.}\,\mathord{.}\,\mathord{.}\,\mathord{.}\,\mathord{.}&0{,}09621\\
\multicolumn{2}{r}{H-K=}&197^{0}\,42'\,23''\\
\multicolumn{2}{r}{H=}&232\quad56\quad42\\
\multicolumn{2}{r}{H-K''=}&221\quad10\quad0.
\end{array}
\]

\SourceLeaf{350}
[Der Ausdruck für die Chorde geht nun in folgenden über, indem man für \(\rho''\) seinen Wert aus VIII einsetzt:
\[
\begin{aligned}
\text{Chorde}^2={}&\Bigl\{\frac{\rho}{M''}h\cos H+g\cos G\\
&\qquad{}+\frac{\cos K''}{M''}\left(\mu-\frac{1}{(r+r'')^3}\right)\Bigr\}^2\\
&+\Bigl\{\frac{\rho}{M''}h\cos J\sin H+\frac{g\sin G}{\cos J}\\
&\qquad{}+\frac{\cos J\sin K''}{M''}\left(\mu-\frac{1}{(r+r'')^3}\right)\Bigr\}^2\\
&+\Bigl\{\frac{\rho}{M''}h\sin J\sin H\\
&\qquad{}+\frac{\sin J\sin K''}{M''}\left(\mu-\frac{1}{(r+r'')^3}\right)\Bigr\}^2\\
={}&\frac{\rho\rho hh}{M''M''}+gg+gg\operatorname{tang}J^2\sin G^2\\
&+\frac{1}{M''M''}\left(\mu-\frac{1}{(r+r'')^3}\right)^2\\
&+\frac{2\rho gh}{M''}\cos(H-G)\\
&+\frac{2\rho h}{M''M''}\left(\mu-\frac{1}{(r+r'')^3}\right)\cos(H-K'')\\
&+\frac{2g}{M''}\left(\mu-\frac{1}{(r+r'')^3}\right)\cos(G-K'')
\end{aligned}
\]
oder, wenn man \(\cos(G-K'')=\cos(H-G)\cos(H-K'')+\sin(H-G)\sin(H-K'')\) setzt und einige leicht ersichtliche Transformationen ausführt:
\[
\begin{aligned}
\text{Chorde}^2={}&\Bigl\{\frac{\rho h}{M''}+g\cos(H-G)\\
&\qquad{}+\frac{\cos(H-K'')}{M''}\left(\mu-\frac{1}{(r+r'')^3}\right)\Bigr\}^2\\
&+(f\sin J\sin(F-N))^2\\
&+\Bigl\{g\sin(H-G)\\
&\qquad{}+\frac{\sin(H-K'')}{M''}\left(\mu-\frac{1}{(r+r'')^3}\right)\Bigr\}^2\mathpunct{\cdot}\Bigr]
\end{aligned}
\]
Hienächst hat man folgende Ausdrücke [wenn man
\[
\begin{aligned}
-u={}&\frac{\rho h}{M''}+g\cos(H-G)\\
&+\frac{\cos(H-K'')}{M''}\left(\mu-\frac{1}{(r+r'')^3}\right)\\
={}&-\frac{\rho''h}{M}+g\cos(H-G)\\
&-\frac{\cos(H-K)}{M}\left(\mu-\frac{1}{(r+r'')^3}\right)
\end{aligned}
\]
setzt, und beachtet, dass \(\dfrac{\sin(H-K'')}{M''}=-\dfrac{\sin(K''-K)}{h}\) ist]:
\[
\begin{array}{l@{\quad}l}
[1]&\begin{aligned}
\rho={}&-\frac{M''}{h}\Bigl\{u+g\cos(H-G)+\frac{\mu\cos(H-K'')}{M''}\\
&\qquad{}-\frac{\cos(H-K'')}{M''}\frac{1}{(r+r'')^3}\Bigr\}
\end{aligned}
\end{array}
\]
\[
\begin{array}{l@{\quad}l}
[2]&\begin{aligned}
\rho''={}&\frac{M}{h}\Bigl\{u+g\cos(H-G)-\frac{\mu\cos(H-K)}{M}\\
&\qquad{}+\frac{\cos(H-K)}{M}\frac{1}{(r+r'')^3}\Bigr\}
\end{aligned}
\end{array}
\]
\[
\begin{array}{l@{\quad}l}
[3]&\begin{aligned}
rr={}&AA+\frac{M''M''}{hh}\Bigl\{u-\frac{\cos(H-K'')}{M''}\frac{1}{(r+r'')^3}\\
&\qquad{}+g\cos(H-G)+\frac{\mu\cos(H-K'')}{M''}-\frac{hR\cos\psi}{M''}\Bigr\}^2
\end{aligned}
\end{array}
\]
\[
\begin{array}{l@{\quad}l}
[4]&\begin{aligned}
r''r''={}&A''A''+\frac{MM}{hh}\Bigl\{u+\frac{\cos(H-K)}{M}\frac{1}{(r+r'')^3}\\
&\qquad{}+g\cos(H-G)-\frac{\mu\cos(H-K)}{M}-\frac{hR''\cos\psi''}{M}\Bigr\}^2
\end{aligned}
\end{array}
\]
\[
\begin{array}{l@{\quad}l}
[5]&\begin{aligned}
\text{Chorde}^2={}&uu+\{f\sin J\sin(F-N)\}^2\\
&+\Bigl\{\frac{\sin(K''-K)}{h}\frac{1}{(r+r'')^3}\\
&\qquad{}+g\sin(H-G)-\frac{\sin(K''-K)}{h}\mu\Bigr\}^2
\end{aligned}
\end{array}
\]

\SourceLeaf{351}
[Numerisch gestalten sich die Gleichungen 1—5 folgendermaassen:]
\[
\begin{array}{l@{\;}l@{\;}r}
-M''&\mathord{.}\,\mathord{.}\,\mathord{.}\,\mathord{.}\,\mathord{.}\,\mathord{.}\,\mathord{.}\,\mathord{.}&0{,}31459\\
h&\mathord{.}\,\mathord{.}\,\mathord{.}\,\mathord{.}\,\mathord{.}\,\mathord{.}\,\mathord{.}\,\mathord{.}\,\mathord{.}\,\mathord{.}\,\mathord{.}&0{,}09621\\\hline
&&0{,}21838
\end{array}
\]
\[
\begin{array}{l@{\;}l@{\;}r}
\cos(H-K'')&\mathord{.}\,\mathord{.}&9{,}87668_n\\
-M''&\mathord{.}\,\mathord{.}\,\mathord{.}\,\mathord{.}\,\mathord{.}\,\mathord{.}\,\mathord{.}\,\mathord{.}&0{,}31459\\\hline
&&9{,}56209_n\\
-\mu&\mathord{.}\,\mathord{.}\,\mathord{.}\,\mathord{.}\,\mathord{.}\,\mathord{.}\,\mathord{.}\,\mathord{.}\,\mathord{.}&9{,}09166_n\\\hline
[\mathrm{I}]&\mathord{.}\,\mathord{.}\,\mathord{.}\,\mathord{.}\,\mathord{.}\,\mathord{.}\,\mathord{.}\,\mathord{.}\,\mathord{.}\,\mathord{.}&8{,}65375
\end{array}
\]
\[
\begin{array}{l@{\;}l@{\;}r}
[M]&\mathord{.}\,\mathord{.}\,\mathord{.}\,\mathord{.}\,\mathord{.}\,\mathord{.}\,\mathord{.}\,\mathord{.}\,\mathord{.}&9{,}97927\\
&\mathord{.}\,\mathord{.}\,\mathord{.}\,\mathord{.}\,\mathord{.}\,\mathord{.}\,\mathord{.}\,\mathord{.}\,\mathord{.}\,\mathord{.}\,\mathord{.}\,\mathord{.}&0{,}09621\\\hline
&&9{,}88306
\end{array}
\]
\[
\begin{array}{l@{\;}l@{\;}r}
[\cos(H-K)]&\mathord{.}&9{,}97892_n\\
{}[M]&\mathord{.}\,\mathord{.}\,\mathord{.}\,\mathord{.}\,\mathord{.}\,\mathord{.}\,\mathord{.}\,\mathord{.}\,\mathord{.}&9{,}97927\\\hline
&&9{,}99965_n\\
&\mathord{.}\,\mathord{.}\,\mathord{.}\,\mathord{.}\,\mathord{.}\,\mathord{.}\,\mathord{.}\,\mathord{.}\,\mathord{.}\,\mathord{.}\,\mathord{.}\,\mathord{.}&9{,}09166_n\\\hline
[\mathrm{II}]&\mathord{.}\,\mathord{.}\,\mathord{.}\,\mathord{.}\,\mathord{.}\,\mathord{.}\,\mathord{.}\,\mathord{.}\,\mathord{.}&9{,}09131
\end{array}
\]
\[
\begin{array}{l@{\;}l@{\;}r}
R\cos\psi&\mathord{.}\,\mathord{.}\,\mathord{.}\,\mathord{.}\,\mathord{.}\,\mathord{.}&9{,}39638_n\\
\dfrac{h}{M''}&\mathord{.}\,\mathord{.}\,\mathord{.}\,\mathord{.}\,\mathord{.}\,\mathord{.}\,\mathord{.}\,\mathord{.}\,\mathord{.}&9{,}78162_n\\[2pt]\hline
&&9{,}17800
\end{array}
\]
\[
\begin{array}{l@{\;}l@{\;}r}
-R''\cos\psi''&\mathord{.}\,\mathord{.}\,\mathord{.}&9{,}84346\\
\left[\dfrac{h}{M}\right]&\mathord{.}\,\mathord{.}\,\mathord{.}\,\mathord{.}\,\mathord{.}\,\mathord{.}\,\mathord{.}\,\mathord{.}\,\mathord{.}&0{,}11694\\[2pt]\hline
&&9{,}96040
\end{array}
\]
\[
\begin{array}{l@{\;}l@{\;}r}
g&\mathord{.}\,\mathord{.}\,\mathord{.}\,\mathord{.}\,\mathord{.}\,\mathord{.}\,\mathord{.}\,\mathord{.}\,\mathord{.}&9{,}29217\\
\cos(H-G)&&9{,}95407\\\hline
[\mathrm{III}]&\mathord{.}\,\mathord{.}\,\mathord{.}\,\mathord{.}\,\mathord{.}\,\mathord{.}&9{,}24624
\end{array}
\]
\[
\begin{array}{l@{\;}l@{\;}r}
[\mathrm{I}+\mathrm{III}]&\mathord{.}\,\mathord{.}\,\mathord{.}\,\mathord{.}&9{,}34508\\
{}[\mathrm{II}+\mathrm{III}]&\mathord{.}\,\mathord{.}\,\mathord{.}&9{,}47668
\end{array}
\]
\[
\begin{array}{l@{\;}l@{\;}r}
f\sin(F-N)&\mathord{.}\,\mathord{.}&9{,}21701_n\\
\sin J&\mathord{.}\,\mathord{.}\,\mathord{.}\,\mathord{.}\,\mathord{.}\,\mathord{.}\,\mathord{.}\,\mathord{.}\,\mathord{.}&9{,}92487\\\hline
&&9{,}14188_n
\end{array}
\]
\[
\begin{array}{l@{\;}l@{\;}r}
\sin(K''-K)&\mathord{.}\,\mathord{.}&9{,}60001_n\\
h&\mathord{.}\,\mathord{.}\,\mathord{.}\,\mathord{.}\,\mathord{.}\,\mathord{.}\,\mathord{.}\,\mathord{.}\,\mathord{.}\,\mathord{.}\,\mathord{.}&0{,}09621\\\hline
&&9{,}50380_n\\
-\mu&\mathord{.}\,\mathord{.}\,\mathord{.}\,\mathord{.}\,\mathord{.}\,\mathord{.}\,\mathord{.}\,\mathord{.}\,\mathord{.}&9{,}09166_n\\\hline
&&8{,}59546
\end{array}
\]
\[
\begin{array}{l@{\;}l@{\;}r}
g&\mathord{.}\,\mathord{.}\,\mathord{.}\,\mathord{.}\,\mathord{.}\,\mathord{.}\,\mathord{.}\,\mathord{.}\,\mathord{.}\,\mathord{.}\,\mathord{.}&9{,}29217\\
\sin(H-G)&\mathord{.}\,\mathord{.}\,\mathord{.}&9{,}64010\\\hline
&&8{,}93227
\end{array}
\]
\[
\begin{array}{l@{\quad}l}
1)&\begin{aligned}\rho={}&(0{,}21838)\Bigl\{u+(9{,}56209_n)\frac{1}{(r+r'')^3}\\
&\qquad{}+(9{,}34508)\Bigr\}\end{aligned}\\[7pt]
2)&\begin{aligned}\rho''={}&(9{,}88306)\Bigl\{u+(9{,}99965_n)\frac{1}{(r+r'')^3}\\
&\qquad{}+(9{,}47668)\Bigr\}\end{aligned}
\end{array}
\]
\[
\begin{array}{l@{\quad}l}
3)&\begin{aligned}rr={}&(9{,}98706)^2+(0{,}21838)^2\Bigl\{u+\frac{(9{,}56209_n)}{(r+r'')^3}\\
&\qquad{}+(9{,}57055)\Bigr\}^2\end{aligned}\\[7pt]
4)&\begin{aligned}r''r''={}&(9{,}86038)^2+(9{,}88306)^2\Bigl\{u+\frac{(9{,}99965_n)}{(r+r'')^3}\\
&\qquad{}+(0{,}08370)\Bigr\}^2\end{aligned}
\end{array}
\]
\[
\begin{array}{l@{\quad}l}
5)&\begin{aligned}\text{Chorde}^2={}&uu+(9{,}14188)^2\\
&+\left\{\frac{(9{,}50380_n)}{(r+r'')^3}+(9{,}09676)\right\}^2.
\end{aligned}
\end{array}
\]

Hiemit sind nun noch folgende Formeln zu verbinden
\[
\begin{array}{l@{\quad}l}
6)&\displaystyle r+r''=\left(\frac{t''-t}{m\BareRoot{\dfrac{108}{5}}}\right)^{\frac{2}{3}}\left(\frac{\BareRoot[5]{\sec W}}{\operatorname{tang}W}\right)^{\frac{2}{3}}
\end{array}
\]
\[
\begin{array}{l@{\quad}l}
7)&\displaystyle\text{Chorde}=\left(\frac{t''-t}{m\BareRoot[4]{\dfrac{135}{4}}}\right)^{\frac{2}{3}}\operatorname{tang}W^{\frac{1}{3}}\mathord{\raisebox{-.6ex}{.}}\sec W^{\frac{2}{15}}.
\end{array}
\]

[Entwickelt man nemlich die bekannte Gleichung, vgl. Band VI, Seite 33,
\SourceLeaf{352}
\[
(r+r''+\text{Chorde})^{\frac{3}{2}}-(r+r''-\text{Chorde})^{\frac{3}{2}}=\frac{t''-t}{m}
\]
nach Potenzen der Grösse \(\dfrac{\text{Chorde}}{r+r''}\mathpunct{\raisebox{.45ex}{,}}\) so wird
\[
\begin{aligned}
\frac{t''-t}{3m}={}&(r+r'')^{\frac{3}{2}}\frac{\text{Chorde}}{r+r''}\Bigl\{1-\frac{1}{24}\left(\frac{\text{Chorde}}{r+r''}\right)^2\\
&\qquad{}-\frac{1}{8\mathord{.}16}\left(\frac{\text{Chorde}}{r+r''}\right)^4\\
&\qquad{}-\frac{3}{16\mathord{.}64}\left(\frac{\text{Chorde}}{r+r''}\right)^6-\mathord{.}\,\mathord{.}\,\mathord{.}\Bigr\}
\end{aligned}
\]
und setzt man
\[
\frac{\text{Chorde}}{r+r''}=\BareRoot{\frac{12}{5}}\mathord{.}\operatorname{tang}W,
\]
so ist
\[
\begin{aligned}
\frac{t''-t}{m}={}&(r+r'')^{\frac{3}{2}}\BareRoot{\frac{108}{5}}\operatorname{tang}W\Bigl\{1-\frac{1}{10}\operatorname{tang}W^2\\
&\qquad{}-\frac{9}{200}\operatorname{tang}W^4-\mathord{.}\,\mathord{.}\,\mathord{.}\Bigr\};
\end{aligned}
\]
nun ist aber
\[
\BareRoot[5]{\cos W}=1-\frac{1}{10}\operatorname{tang}W^2+\frac{11}{200}\operatorname{tang}W^4-\mathord{.}\,\mathord{.}\,\mathord{.}\Bigr\},
\]
also genähert
\[
\frac{t''-t}{m\BareRoot{\dfrac{108}{5}}}=(r+r'')^{\frac{3}{2}}\operatorname{tang}W\BareRoot[5]{\cos W},
\]
womit die Gleichungen 6 und 7 sich ergeben.]
\[
\begin{array}{l@{\;}l@{\;}r}
m&\mathord{.}\,\mathord{.}\,\mathord{.}\,\mathord{.}\,\mathord{.}\,\mathord{.}\,\mathord{.}\,\mathord{.}\,\mathord{.}&0{,}986\,2673\\
\BareRoot{\dfrac{108}{5}}&\mathord{.}\,\mathord{.}\,\mathord{.}\,\mathord{.}\,\mathord{.}&0{,}667\,2269\\[3pt]\hline
&&1{,}653\,4942\\
t''-t&\mathord{.}\,\mathord{.}\,\mathord{.}\,\mathord{.}\,\mathord{.}\,\mathord{.}&1{,}147\,6544\\\hline
&&0{,}505\,8398
\end{array}
\]
\[
\begin{array}{l@{\;}l@{\;}r}
&\mathord{.}\,\mathord{.}\,\mathord{.}\,\mathord{.}\,\mathord{.}\,\mathord{.}\,\mathord{.}\,\mathord{.}\,\mathord{.}\,\mathord{.}&0{,}986\,2673\\
\BareRoot[4]{\dfrac{135}{4}}&\mathord{.}\,\mathord{.}\,\mathord{.}\,\mathord{.}\,\mathord{.}&0{,}382\,0684\\[3pt]\hline
&&1{,}368\,3357\\
&\mathord{.}\,\mathord{.}\,\mathord{.}\,\mathord{.}\,\mathord{.}\,\mathord{.}\,\mathord{.}\,\mathord{.}\,\mathord{.}\,\mathord{.}&1{,}147\,6544\\\hline
&&0{,}220\,6813
\end{array}
\]
\[
\begin{array}{l@{\quad}l}
[6]&\displaystyle r+r''=\frac{(9{,}66277)}{\left(\dfrac{\operatorname{tang}W}{\BareRoot[5]{\sec W}}\right)^{\frac{2}{3}}}\mathpunct{\raisebox{.45ex}{,}}\\[8pt]
[7]&\displaystyle\text{Chorde}=(9{,}85288)\operatorname{tang}W^{\frac{1}{3}}\mathord{.}\sec W^{\frac{2}{15}}.
\end{array}
\]

Nach einem vorausgesetzten Werthe von \(W\) berechnet man aus den Formeln 6 und 7 die Chorde und \(r+r''\); demnächst aus der Formel 5 den Werth von \(u\) und endlich aus 3 und 4 den zweiten Werth von \(r+r''\).

Noch zweckmässiger ist folgende Formel
\[
\begin{array}{l@{\quad}l}
[8]&\begin{aligned}
\text{Chorde}&=\frac{\dfrac{t''-t}{3m}}{\BareRoot{(r+r'')}\mathord{.}\BareRoot[14]{\left(1-\dfrac{(t''-t)^2}{\dfrac{108}{7}mm(r+r'')^3}\right)}}\\[5pt]
&=\frac{(9{,}68427)}{\BareRoot{(r+r'')}\mathord{.}\BareRoot[14]{\left(1-\dfrac{(9{,}1344484)}{(r+r'')^3}\right)}}
\end{aligned}
\end{array}
\]
\SourceLeaf{353}
[aus welcher man mit einem vorausgesetzten Werth von \(r+r''\) die Chorde berechnet, worauf, wie vorher, die Gleichung 5 den Werth von \(u\) und 3 und 4 den zweiten Werth von \(r+r''\) liefern.

Man findet die Gleichung 8, indem man die obige Reihe umkehrt, womit
\[
\begin{aligned}
\text{Chorde}={}&\frac{t''-t}{3m\BareRoot{(r+r'')}}\Bigl\{1+\frac{1}{2\mathord{.}108}\left(\frac{t''-t}{m(r+r'')^{\frac{3}{2}}}\right)^2\\
&\qquad{}+\frac{15}{8\mathord{.}108^2}\left(\frac{t''-t}{m(r+r'')^{\frac{3}{2}}}\right)^4\\
&\qquad{}+\frac{177}{16\mathord{.}108^3}\left(\frac{t''-t}{m(r+r'')^{\frac{3}{2}}}\right)^6+\mathord{.}\,\mathord{.}\,\mathord{.}\Bigr\}
\end{aligned}
\]
wird; da nun
\[
\begin{aligned}
\frac{1}{\BareRoot[14]{\left(1-\dfrac{(t''-t)^2}{\dfrac{108}{7}mm(r+r'')^3}\right)}}
&=1+\frac{1}{2\mathord{.}108}\left(\frac{t''-t}{m(r+r'')^{\frac{3}{2}}}\right)^2\\
&\quad{}+\frac{15}{8\mathord{.}108^2}\left(\frac{t''-t}{m(r+r'')^{\frac{3}{2}}}\right)^4\\
&\quad{}+\frac{145}{16\mathord{.}108^3}\left(\frac{t''-t}{m(r+r'')^{\frac{3}{2}}}\right)^6+\mathord{.}\,\mathord{.}\,\mathord{.},
\end{aligned}
\]
so folgt die Gleichung 8, welche eine noch grössere Annäherung gibt, als die Gleichungen 6 und 7].
\begin{center}Numerische Rechnung.\end{center}
Vorausgesetzter Werth \(\log(r+r'')=0{,}42622\).
\[
\begin{array}{l@{\;}l@{\;}r}
(r+r'')^3&\mathord{.}\,\mathord{.}&1{,}27866\\
&\mathord{.}\,\mathord{.}\,\mathord{.}\,\mathord{.}\,\mathord{.}\,\mathord{.}\,\mathord{.}&9{,}13445\\\hline
&&[7{,}85579]
\end{array}
\]
\[
\begin{array}{l@{\;}l@{\;}r}
&&9{,}50380_n\\
(r+r'')^3&\mathord{.}\,\mathord{.}&1{,}27866\\\hline
&&8{,}22514_n\\
&\mathord{.}\,\mathord{.}\,\mathord{.}\,\mathord{.}\,\mathord{.}\,\mathord{.}\,\mathord{.}\,\mathord{.}&9{,}09676\\\hline
&&9{,}03408
\end{array}
\]
\[
\begin{array}{l@{\;}l@{\;}r}
&&9{,}56209_n\\
(r+r'')^3&\mathord{.}&1{,}27866\\\hline
&&8{,}28343_n\\
&&9{,}57055\\
u&\mathord{.}\,\mathord{.}\,\mathord{.}\,\mathord{.}\,\mathord{.}\,\mathord{.}&9{,}37691\\\hline
&&9{,}77157
\end{array}
\]
\[
\begin{array}{l@{\;}l@{\;}r}
&&9{,}99965_n\\
&\mathord{.}\,\mathord{.}\,\mathord{.}\,\mathord{.}\,\mathord{.}&1{,}27866\\\hline
&&8{,}72099_n\\
&&0{,}08370\\
&\mathord{.}\,\mathord{.}\,\mathord{.}\,\mathord{.}\,\mathord{.}&9{,}37691\\\hline
&&0{,}14555
\end{array}
\]
\[
\begin{array}{l@{\;}l@{\;}r}
\BareRoot{(r+r'')}&\mathord{.}&0{,}21311\\
&\mathord{.}\,\mathord{.}\,\mathord{.}\,\mathord{.}\,\mathord{.}&-0{,}00022\\\hline
[\text{Nenner}]&\mathord{.}&0{,}21289\\
&\mathord{.}\,\mathord{.}\,\mathord{.}\,\mathord{.}\,\mathord{.}\,\mathord{.}\,\mathord{.}&9{,}68427\\\hline
\text{Chorde}&\mathord{.}\,\mathord{.}&9{,}47138
\end{array}
\]
\[
\begin{array}{l@{\;}l@{\;}r}
&&[8{,}06816]\\
&&[8{,}28376]\\
\multicolumn{3}{r}{[\text{Chorde}^2\,\mathord{.}\quad8{,}94276]}\\\hline
u&\mathord{.}\,\mathord{.}\,\mathord{.}\,\mathord{.}\,\mathord{.}\,\mathord{.}&9{,}37691
\end{array}
\]
\[
\begin{array}{l@{\;}l@{\;}r}
&&\left[\begin{array}{@{}r@{}}9{,}54314\\0{,}43676\\\hline9{,}97990\\9{,}97412\end{array}\right]\\\hline
r&\mathord{.}\,\mathord{.}\,\mathord{.}\,\mathord{.}\,\mathord{.}\,\mathord{.}&0{,}13902
\end{array}
\]
\[
\begin{array}{l@{\;}l@{\;}r}
&&\left[\begin{array}{@{}r@{}}0{,}29110\\9{,}76612\\\hline0{,}05722\\9{,}72076\end{array}\right]\\\hline
r''&\mathord{.}\,\mathord{.}\,\mathord{.}&0{,}11091
\end{array}
\]
\noindent wie vorausgesetzt, wird \(r+r''\quad\mathord{.}\,\mathord{.}\,\mathord{.}\quad0{,}42622\).
\SourceLeaf{354}
\[
\begin{array}{l@{\;}l@{\;}r}
&&8{,}28343_n\\
&&9{,}34508\\
u&\mathord{.}\,\mathord{.}\,\mathord{.}\,\mathord{.}\,\mathord{.}&9{,}37691\\\hline
&&9{,}64378\\
&&0{,}21838\\\hline
\rho&\mathord{.}\,\mathord{.}\,\mathord{.}\,\mathord{.}\,\mathord{.}&9{,}86216
\end{array}
\]
\[
\begin{array}{l@{\;}l@{\;}r}
&&8{,}72099_n\\
&&9{,}47668\\
&\mathord{.}\,\mathord{.}\,\mathord{.}\,\mathord{.}\,\mathord{.}\,\mathord{.}&9{,}37691\\\hline
&&9{,}68599\\
&&9{,}88306\\\hline
\rho''&\mathord{.}\,\mathord{.}\,\mathord{.}\,\mathord{.}&9{,}56905
\end{array}
\]
\[
\begin{array}{l@{\;}l@{\;}r}
\rho&\mathord{.}\,\mathord{.}\,\mathord{.}\,\mathord{.}\,\mathord{.}\,\mathord{.}\,\mathord{.}&9{,}86216\\
\cos\beta&\mathord{.}\,\mathord{.}\,\mathord{.}\,\mathord{.}&9{,}94168\\\hline
\rho\cos\beta&\mathord{.}\,\mathord{.}\,\mathord{.}&9{,}80334
\end{array}
\]
\[
\begin{array}{l@{\;}l@{\;}r}
\rho''&\mathord{.}\,\mathord{.}\,\mathord{.}\,\mathord{.}\,\mathord{.}\,\mathord{.}\,\mathord{.}&9{,}56905\\
\cos\beta''&\mathord{.}\,\mathord{.}\,\mathord{.}&9{,}99350\\\hline
\rho''\cos\beta''&\mathord{.}\,\mathord{.}&9{,}56255
\end{array}
\]

Wir fügen zugleich die vollständige Berechnung der Elemente bei [vgl. Band VI, Seite 34—36]:
\[
\alpha-\odot=253^{0}\,28'\,57''
\]
\[
\begin{array}{l@{\;}l@{\;}r}
\rho\cos\beta&\mathord{.}\,\mathord{.}\,\mathord{.}\,\mathord{.}\,\mathord{.}\,\mathord{.}\,\mathord{.}\,\mathord{.}\,\mathord{.}\,\mathord{.}\,\mathord{.}\,\mathord{.}\,\mathord{.}\,\mathord{.}&9{,}80384\\
\cos(\alpha-\odot)&\mathord{.}\,\mathord{.}\,\mathord{.}\,\mathord{.}\,\mathord{.}\,\mathord{.}\,\mathord{.}\,\mathord{.}\,\mathord{.}\,\mathord{.}&9{,}45379_n\\
\sin(\alpha-\odot)&\mathord{.}\,\mathord{.}\,\mathord{.}\,\mathord{.}\,\mathord{.}\,\mathord{.}\,\mathord{.}\,\mathord{.}\,\mathord{.}\,\mathord{.}&9{,}98170_n\\
\operatorname{tang}\beta&\mathord{.}\,\mathord{.}\,\mathord{.}\,\mathord{.}\,\mathord{.}\,\mathord{.}\,\mathord{.}\,\mathord{.}\,\mathord{.}\,\mathord{.}\,\mathord{.}\,\mathord{.}\,\mathord{.}\,\mathord{.}&9{,}74435\\\hline
\end{array}
\]
\[
[\alpha''-\odot''=]\ 225^{0}\,16'\,43''
\]
\[
\begin{array}{l@{\;}l@{\;}r}
\left[\begin{array}{@{}l@{}}\rho''\cos\beta''\quad\mathord{.}\,\mathord{.}\,\mathord{.}\,\mathord{.}\\\cos(\alpha''-\odot'')\\\sin(\alpha''-\odot'')\\\operatorname{tang}\beta''\quad\mathord{.}\,\mathord{.}\,\mathord{.}\,\mathord{.}\,\mathord{.}\end{array}\right]&
\begin{array}{@{}l@{}}\mathord{.}\,\mathord{.}\,\mathord{.}\,\mathord{.}\,\mathord{.}\\\mathord{.}\,\mathord{.}\,\mathord{.}\,\mathord{.}\,\mathord{.}\\\mathord{.}\quad\mathord{.}\,\mathord{.}\,\mathord{.}\\\mathord{.}\,\mathord{.}\,\mathord{.}\,\mathord{.}\,\mathord{.}\end{array}&
\begin{array}{@{}r@{}}9{,}56255\\9{,}84736_n\\9{,}85159_n\\9{,}24127\end{array}\\\hline
\end{array}
\]
\[
\begin{array}{l@{\;}l@{\;}r}
[\rho\cos\beta\cos(\alpha-\odot)]&\mathord{.}\,\mathord{.}\,\mathord{.}\,\mathord{.}&9{,}25763_n\\
-R&\mathord{.}\,\mathord{.}\,\mathord{.}\,\mathord{.}\,\mathord{.}\,\mathord{.}\,\mathord{.}\,\mathord{.}\,\mathord{.}\,\mathord{.}\,\mathord{.}\,\mathord{.}\,\mathord{.}\,\mathord{.}\,\mathord{.}&0{,}00091_n\\\hline
\left[\begin{array}{@{}l@{}}r\cos b\cos(\lambda-\odot)\\r\cos b\sin(\lambda-\odot)\end{array}\right]&
\begin{array}{@{}l@{}}\mathord{.}\,\mathord{.}\\\mathord{.}\,\mathord{.}\end{array}&
\begin{array}{@{}r@{}}0{,}07301_n\\9{,}78554_n\end{array}\\\hline
\multicolumn{3}{l}{\left[\begin{array}{l@{\;}l@{\;}r}
\cos(\lambda-\odot)&\mathord{.}\,\mathord{.}\,\mathord{.}\,\mathord{.}\,\mathord{.}\,\mathord{.}\,\mathord{.}\,\mathord{.}\,\mathord{.}&9{,}94876_n\\
r\cos b&\mathord{.}\,\mathord{.}\,\mathord{.}\,\mathord{.}\,\mathord{.}\,\mathord{.}\,\mathord{.}\,\mathord{.}\,\mathord{.}\,\mathord{.}\,\mathord{.}\,\mathord{.}\,\mathord{.}&0{,}12425\\
r\sin b&\mathord{.}\,\mathord{.}\,\mathord{.}\,\mathord{.}\,\mathord{.}\,\mathord{.}\,\mathord{.}\,\mathord{.}\,\mathord{.}\,\mathord{.}\,\mathord{.}\,\mathord{.}\,\mathord{.}&9{,}54819
\end{array}\right]}\\\hline
\cos b&\mathord{.}\,\mathord{.}\,\mathord{.}\,\mathord{.}\,\mathord{.}\,\mathord{.}\,\mathord{.}\,\mathord{.}\,\mathord{.}\,\mathord{.}\,\mathord{.}\,\mathord{.}\,\mathord{.}\,\mathord{.}\,\mathord{.}&9{,}98522\\
r&\mathord{.}\,\mathord{.}\,\mathord{.}\,\mathord{.}\,\mathord{.}\,\mathord{.}\,\mathord{.}\,\mathord{.}\,\mathord{.}\,\mathord{.}\,\mathord{.}\,\mathord{.}\,\mathord{.}\,\mathord{.}\,\mathord{.}\,\mathord{.}\,\mathord{.}\,\mathord{.}&0{,}13903
\end{array}
\]
\[
\begin{array}{l@{\;}l@{\;}r}
[\rho''\cos\beta''\cos(\alpha''-\odot'')]&\mathord{.}&9{,}40991_n\\
{}[-R'']&\mathord{.}\,\mathord{.}\,\mathord{.}\,\mathord{.}\,\mathord{.}\,\mathord{.}\,\mathord{.}\,\mathord{.}\,\mathord{.}\,\mathord{.}\,\mathord{.}\,\mathord{.}\,\mathord{.}\,\mathord{.}&0{,}00260_n\\\hline
\left[\begin{array}{@{}l@{}}r''\cos b''\cos(\lambda''-\odot'')\\r''\cos b''\sin(\lambda''-\odot'')\end{array}\right]&&
\begin{array}{@{}r@{}}0{,}10140_n\\9{,}41414_n\end{array}\\\hline
\multicolumn{3}{l}{\left[\begin{array}{l@{\;}l@{\;}r}
\cos(\lambda''-\odot'')&\mathord{.}\,\mathord{.}\,\mathord{.}\,\mathord{.}\,\mathord{.}\,\mathord{.}\,\mathord{.}&9{,}99102_n\\
r''\cos b''&\mathord{.}\,\mathord{.}\,\mathord{.}\,\mathord{.}\,\mathord{.}\,\mathord{.}\,\mathord{.}\,\mathord{.}\,\mathord{.}\,\mathord{.}\,\mathord{.}\,\mathord{.}&0{,}11038\\
r''\sin b''&\mathord{.}\,\mathord{.}\,\mathord{.}\,\mathord{.}\,\mathord{.}\,\mathord{.}\,\mathord{.}\,\mathord{.}\,\mathord{.}\,\mathord{.}\,\mathord{.}\,\mathord{.}&8{,}80382
\end{array}\right]}\\\hline
\left[\begin{array}{@{}l@{}}\cos b''\\r''\quad\mathord{.}\,\mathord{.}\,\mathord{.}\end{array}\right]&
\begin{array}{@{}l@{}}\mathord{.}\,\mathord{.}\,\mathord{.}\,\mathord{.}\,\mathord{.}\,\mathord{.}\,\mathord{.}\,\mathord{.}\,\mathord{.}\,\mathord{.}\,\mathord{.}\,\mathord{.}\\\mathord{.}\,\mathord{.}\,\mathord{.}\,\mathord{.}\,\mathord{.}\,\mathord{.}\,\mathord{.}\,\mathord{.}\,\mathord{.}\,\mathord{.}\,\mathord{.}\,\mathord{.}\end{array}&
\begin{array}{@{}r@{}}9{,}99947\\0{,}11091\end{array}
\end{array}
\]
\[
\begin{array}{r@{\;}l}
[\lambda-\odot]=&207^{0}\,17'\,14''\\
\odot=&17\quad47\quad41\\
\lambda=&225\quad4\quad55\\
b=&14\quad51\quad54
\end{array}
\]
\[
\begin{array}{r@{\;}l}
[\lambda''-\odot'']=&191^{0}\,36'\,39''\\
\left[\begin{array}{@{}r@{}}\odot''=\\\lambda''=\\b''=\end{array}\right]&
\begin{array}{@{}r@{\quad}r@{\quad}r@{}}31&31&25\\223&8&4\\2&49&34\end{array}
\end{array}
\]
\SourceLeaf{355}
\[
\begin{array}{l@{\;}l@{\;}r}
-\operatorname{tang}b&\mathord{.}\,\mathord{.}\,\mathord{.}\,\mathord{.}\,\mathord{.}\,\mathord{.}\,\mathord{.}&9{,}42394_n\\
\cos(\lambda''-\lambda)&\mathord{.}\,\mathord{.}\,\mathord{.}\,\mathord{.}\,\mathord{.}\,\mathord{.}&9{,}99975\\\hline
&&9{,}42369_n\\
{}[\operatorname{tang}b'']&\mathord{.}\,\mathord{.}\,\mathord{.}\,\mathord{.}\,\mathord{.}\,\mathord{.}\,\mathord{.}\,\mathord{.}&8{,}69344\\\hline
&&9{,}33426_n\\
\sin(\lambda''-\lambda)&\mathord{.}\,\mathord{.}\,\mathord{.}\,\mathord{.}\,\mathord{.}\,\mathord{.}&8{,}53127_n\\\hline
\end{array}
\]
\[
\begin{array}{l@{\;}l@{\;}r}
\left[\begin{array}{@{}l@{}}\operatorname{tang}i\sin(\lambda-\AscendingNode)\\\operatorname{tang}i\cos(\lambda-\AscendingNode)\end{array}\right]&&
\begin{array}{@{}r@{}}9{,}42394_n\\0{,}80299_n\end{array}\\\hline
[\cos(\lambda-\AscendingNode)]&\mathord{.}\,\mathord{.}\,\mathord{.}\,\mathord{.}\,\mathord{.}&9{,}99962_n
\end{array}
\]
\[
\begin{array}{r@{\;}l}
\lambda-\AscendingNode=&182^{0}\,23'\,32''\\
\lambda=&225\quad4\quad55\\\hline
\AscendingNode=&42\quad41\quad23\\
\lambda''-\AscendingNode=&180\quad26\quad41\\
\operatorname{tang}i\quad\mathord{.}\,\mathord{.}\,\mathord{.}\,\mathord{.}\,\mathord{.}\,\mathord{.}\,\mathord{.}\,\mathord{.}\,\mathord{.}\,\mathord{.}&0{,}80337\\
i=&81^{0}\,3'\,45''
\end{array}
\]
\[
\begin{array}{l@{\;}l@{\;}r}
\operatorname{tang}(\lambda-\AscendingNode)&\mathord{.}\,\mathord{.}\,\mathord{.}\,\mathord{.}\,\mathord{.}&8{,}62095_n\\
\operatorname{tang}(\lambda''-\AscendingNode)&\mathord{.}\,\mathord{.}\,\mathord{.}\,\mathord{.}&7{,}88997_n\\
{}[\cos i]&\mathord{.}\,\mathord{.}\,\mathord{.}\,\mathord{.}\,\mathord{.}\,\mathord{.}\,\mathord{.}\,\mathord{.}\,\mathord{.}\,\mathord{.}&9{,}19132\\\hline
\end{array}
\]
\[
\begin{array}{r@{\;}l}
v-\AscendingNode=&195^{0}\,3'\,8''\\
v''-\AscendingNode=&182\quad51\quad37\\
\AscendingNode=&42\quad41\quad23\\\hline
v=&237\quad44\quad31\\
v''=&225\quad33\quad0\\
\frac{1}{2}(v''-v)=&-6\quad5\quad45{,}5
\end{array}
\]
\[
\begin{array}{l@{\;}l@{\;}r}
\dfrac{1}{\BareRoot{r}}&\mathord{.}\,\mathord{.}\,\mathord{.}\,\mathord{.}\,\mathord{.}\,\mathord{.}\,\mathord{.}\,\mathord{.}\,\mathord{.}\,\mathord{.}\,\mathord{.}\,\mathord{.}&9{,}93048\\[3pt]
\operatorname{tang}\frac{1}{2}(v''-v)&\mathord{.}\,\mathord{.}\,\mathord{.}\,\mathord{.}&9{,}02856_n\\\hline
\sin\frac{1}{2}(v''-v)&\mathord{.}\,\mathord{.}\,\mathord{.}\quad\mathord{.}&9{,}02610_n\\
\dfrac{1}{\BareRoot{r''}}&\mathord{.}\,\mathord{.}\,\mathord{.}\,\mathord{.}\,\mathord{.}\,\mathord{.}\,\mathord{.}\,\mathord{.}\,\mathord{.}\,\mathord{.}\,\mathord{.}\,\mathord{.}&9{,}94454\\[3pt]\hline
\end{array}
\]
\[
\begin{array}{l@{\;}r}
&0{,}90192_n\\
&0{,}91844\\\hline
\left[\begin{array}{@{}l@{}}\dfrac{1}{\BareRoot{q}}\sin\frac{1}{2}(v-\Pi)\\[4pt]\dfrac{1}{\BareRoot{q}}\cos\frac{1}{2}(v-\Pi)\end{array}\right]&
\begin{array}{@{}r@{}}9{,}49042\\[4pt]9{,}93048\end{array}\\\hline
\end{array}
\]
\[
\begin{array}{l@{\;}l@{\;}r}
[\cos\frac{1}{2}(v-\Pi)]&\mathord{.}\,\mathord{.}\,\mathord{.}\,\mathord{.}&9{,}97311\\
\multicolumn{2}{r}{\frac{1}{2}(v-\Pi)=}&19^{0}\,57'\,8''\\
q&\mathord{.}\,\mathord{.}\,\mathord{.}\,\mathord{.}\,\mathord{.}\,\mathord{.}\,\mathord{.}\,\mathord{.}\,\mathord{.}\,\mathord{.}\,\mathord{.}\,\mathord{.}\,\mathord{.}&0{,}08526\\
\multicolumn{2}{r}{v-\Pi=}&39^{0}\,54'\,16''\\
\multicolumn{2}{r}{[v=]}&237\quad44\quad31\\\hline
\multicolumn{2}{r}{\Pi=}&197\quad50\quad15\\
\multicolumn{2}{r}{v''-\Pi=}&27\quad42\quad45
\end{array}
\]
\[
\begin{array}{l@{\;}l@{\;}r}
n&\mathord{.}\quad\mathord{.}\,\mathord{.}\,\mathord{.}\,\mathord{.}\,\mathord{.}\,\mathord{.}\,\mathord{.}\,\mathord{.}\,\mathord{.}&0{,}03987\\
q^{\frac{3}{2}}&\mathord{.}\,\mathord{.}\,\mathord{.}\,\mathord{.}\,\mathord{.}\,\mathord{.}\,\mathord{.}\,\mathord{.}\,\mathord{.}\,\mathord{.}\,\mathord{.}&0{,}12789\\\hline
&&0{,}16776
\end{array}
\]
\[
\begin{array}{c@{\qquad}c}
\begin{array}{l@{\;}l@{\;}r}
[M]&\mathord{.}\,\mathord{.}\,\mathord{.}\,\mathord{.}\,\mathord{.}\,\mathord{.}\,\mathord{.}\,\mathord{.}\,\mathord{.}&1{,}45367\\
&&[0{,}16776]\\\hline
\multicolumn{3}{l}{[T-t\quad\mathord{.}\,\mathord{.}\,\mathord{.}\,\mathord{.}\,\mathord{.}\,\mathord{.}\,\mathord{.}\quad1{,}62143]}\\
\left[\begin{array}{@{}r@{}}T-t=\\t=\end{array}\right]&&\begin{array}{@{}r@{}}41{,}8245\\7{,}5500\end{array}\\\hline
\end{array}
&
\begin{array}{l@{\;}l@{\;}r}
[M'']&\mathord{.}\,\mathord{.}\,\mathord{.}\,\mathord{.}\,\mathord{.}\,\mathord{.}\,\mathord{.}\,\mathord{.}&1{,}27589\\
&&[0{,}16776]\\\hline
\multicolumn{3}{l}{[T-t''\quad\mathord{.}\,\mathord{.}\,\mathord{.}\,\mathord{.}\,\mathord{.}\,\mathord{.}\quad1{,}44365]}\\
\left[\begin{array}{@{}r@{}}T-t''=\\t''=\end{array}\right]&&\begin{array}{@{}r@{}}27{,}7746\\21{,}5993\end{array}\\\hline
\end{array}\\
\multicolumn{2}{c}{\underbrace{\begin{array}{c@{\qquad}c}[T=]\ 49{,}3745&[T=]\ 49{,}3739\end{array}}_{\textstyle[\text{Mittel }T=]\ 49{,}3742}}\\[4pt]
\multicolumn{2}{c}{[=\text{ Mai }19{,}3742.]}
\end{array}
\]
\par\medskip\centerline{\rule{.18\linewidth}{.4pt}}\medskip
